\documentclass[12pt,a4paper]{article}
\usepackage[utf8]{inputenc}
\usepackage[spanish]{babel}
\usepackage{graphicx}
\usepackage{booktabs}
\usepackage{tabularx}
\usepackage[hidelinks]{hyperref}
\usepackage{geometry}
\geometry{top=2.5cm, bottom=2.5cm, left=3cm, right=3cm}

\title{\textbf{Escenarios 9 al 12: Visualizaci\'on y Plan de Mejora Continua}}
\author{
\textbf{Asignatura:} Aseguramiento de la Calidad del Software --- NRC 30733 \\
\textbf{Docente:} Ing. Diego Gamboa \\
\textbf{Grupo 7 --- Calidad en Uso} \\
\textbf{Integrantes:} Kevin Amagua\~na, Jairo Bonilla, Darwin Panchez, Eduardo Mortensen
}
\date{\today}

\begin{document}

\maketitle

\section{Descripci\'on General}
Este documento sirve como base para integrar los reportes correspondientes a los Escenarios 9 al 12 del Taller de Calidad de Uso. Estos escenarios completan el ciclo anal\'itico partiendo de los datos autom\'aticamente generados en el Escenario 8.

\section{Escenario 9: Construcci\'on de indicadores (KPI)}
A partir de los archivos exportados (como \texttt{indicadores.json}), se elaboran tableros de control y gr\'aficas visuales (dashboards) que permiten resumir el estado de la calidad en uso de MediSalud HIS en tiempo real.

\section{Escenario 10: Interpretaci\'on de resultados}
Con los indicadores disponibles, se procede al an\'alisis causal y diagn\'ostico. Si el umbral para RNF-01 (registro $\le$ 8 segundos) no se cumple, el equipo investigar\'a la variaci\'on entre horas pico y las diferencias geogr\'aficas para comprender la ra\'iz de las deficiencias de Eficiencia y Satisfacci\'on.

\section{Escenario 11: Presentaci\'on ejecutiva para directivos}
Los hallazgos t\'ecnicos son consolidados y traducidos al lenguaje del negocio, evidenciando el impacto directo de la "baja calidad en uso" sobre el flujo de caja, la efectividad cl\'inica y la experiencia del paciente, justificando las inversiones tecnol\'ogicas necesarias.

\section{Escenario 12: Plan de mejora continua}
Finalmente, se dise\~nan acciones correctivas y de mejora orientadas a optimizar el rendimiento (por ejemplo, optimizar el backend del m\'odulo HCE o el motor de base de datos) y reducir la latencia general, reingresando dichos objetivos al ciclo iterativo del desarrollo.

\section{Conclusiones Finales}
Los escenarios 9 al 12 consolidan el programa, demostrando que ISO/IEC 25022 no es s\'olo un marco conceptual de medici\'on, sino una herramienta gerencial pr\'actica indispensable para dirigir los proyectos de TI basados en evidencia y orientados a la percepci\'on y eficacia de los usuarios finales.

\end{document}
